\section{Ablations}
\label{sec:ablations}
All ablations use the PoE variant, 3 seeds and the same evaluation run within each table. The training-target ablation runs on the general-object benchmarks (T-LESS, SYMSOL-T). The ablations of the head architecture, the orbit scoring rule and the augmentation were run on the spacecraft application (SPEED+), where the synthetic-to-real shift makes their effects largest; we do not assume that their magnitudes carry over to the other benchmarks.

\begin{table*}[t]
  \centering\footnotesize
  \setlength{\tabcolsep}{4pt}
  \resizebox{\linewidth}{!}{\begin{tabular}{@{}l l l l cc c@{}}
\toprule
Component & Benchmark & Matched comparison & Metric & Without & With & Extra cost \\
\midrule
Multi-path targets & T-LESS & A0 vs.\ A2, PoE + orbit & test R@10 (\%) & 73.5 & 79.0 & none \\  % src: EXP-106 RUN-20261027-201558 / RUN-20261028-221625
 & SYMSOL-T & A0 vs.\ A2, DINOv3 + orbit & test $E_R$ & 2.62 & 2.07 & none \\  % src: EXP-106 RUN-20261028-014408 / RUN-20261028-120353
\addlinespace[2pt]
Canonical labels & SYMSOL-T & A0 vs.\ A1, PoE + orbit & val $E_R$ & 2.65 & 9.73 & none \\  % src: EXP-106 (canonical labels break orbit voting)
\addlinespace[2pt]
Orbit voting & T-LESS & A2 PoE, 1 view vs.\ orbit & val $E_R$ & 10.82 & 8.28 & 4.4$\times$ \\  % src: EXP-106 RUN-20261028-221625; EXP-112
 & SYMSOL-T & A2 DINOv3, 1 view vs.\ orbit & val $E_R$ & 2.54 & 2.07 & 4.7$\times$ \\  % src: EXP-106 RUN-20261028-120353
 & YCB-V & DINOv3, 1 view vs.\ orbit & val $E_R$ & 8.97 & 11.53 & 4.7$\times$ \\  % src: EXP-122 RUN-20261027-200852
 & SPEED+ & aug.\ tree PoE, 1 view vs.\ orbit & LB / SL $E_R$ & 17.3 / 20.2 & 12.3 / 12.3 & 4.4$\times$ \\  % src: EXP-109 RUN-20261016-014754
\addlinespace[2pt]
Distillation & SPEED+ & tree vs.\ student, PoE + orbit & LB / SL $E_R$ & 12.30 / 12.25 & 11.21 / 11.09 & none \\  % src: EXP-117 RUN-20261020-114253 (family tree vs tree:teacher_variant=PoE, orbit_product)
 & T-LESS & A2 tree vs.\ student, PoE, 1 view / orbit & val $E_R$ & 10.82 / 8.28 & 10.08 / 7.88 & none \\  % src: EXP-119 RUN-20261101-065144 vs EXP-106 RUN-20261028-221625
 & SYMSOL-T & A2 tree vs.\ student, DINOv3, 1 view / orbit & val $E_R$ & 2.54 / 2.07 & 2.52 / 2.16 & none \\  % src: EXP-119 RUN-20261031-120353 vs EXP-106 RUN-20261028-120353
\addlinespace[2pt]
Geometry fusion (PoE) & SPEED+ & DINOv3 vs.\ PoE, orbit & LB $E_R$ & 17.4 & 12.3 & 4.0$\times$ \\  % src: EXP-113 RUN-20261017-042606; EXP-112
 & T-LESS & DINOv3 vs.\ PoE, A2 orbit & val $E_R$ & 8.61 & 8.28 & 4.0$\times$ \\  % src: EXP-106 RUN-20261028-221625
 & SYMSOL-T & DINOv3 vs.\ PoE, A2 orbit & val $E_R$ & 2.07 & 2.10 & 4.0$\times$ \\  % src: EXP-106 RUN-20261028-120353
 & YCB-V & DINOv3 vs.\ PoE, 1 view & val $E_R$ & 8.97 & 10.52 & 4.4$\times$ \\  % src: EXP-122 RUN-20261027-200852
\addlinespace[2pt]
Synthetic training data & YCB-V & real only vs.\ + renders, val-selected & test $E_R$ & 69.4 & 12.9 & none \\  % src: EXP-064 RUN-20261019-224920; EXP-122
\addlinespace[2pt]
Hierarchical tree & SPEED+ & flat, same 131,072 leaves vs.\ tree, PoE + orbit & synth.\ $E_R$ & 5.34 & 3.87 & $\approx14\times$ fewer params \\  % src: EXP-127 RUN-20261103-143538
 & SPEED+ & flat, balanced Hopf 1771$\times$74 vs.\ tree, PoE + orbit & synth.\ $E_R$ & 5.28 & 3.87 &  \\  % src: EXP-127 RUN-20261103-143538
\addlinespace[2pt]
Translation block & T-LESS & box only vs.\ block (orbit mean) & test $E_T$ (cm) & 6.3 & 2.6 & +0.33\,M \\  % src: EXP-126 RUN-20261030-085948
 & YCB-V & box only vs.\ block (1 view) & test $E_T$ (cm) & 8.1 & 1.6 & +0.33\,M \\  % src: EXP-126 RUN-20261030-044920
\bottomrule
\end{tabular}}
  \caption{\textbf{Component contributions under matched controls.} Each row changes one component and keeps features, recipe, seeds and evaluation fixed; values are means over 3 seeds ($^\circ$ unless stated; lower is better except recall). Extra cost: end-to-end inference time relative to the configuration without the component, measured on SPEED+ crops (\cref{tab:cost_infer}; the roll-homography warp used for T-LESS and YCB-V is not timed), or added parameters. Source runs are recorded next to each row in the table generator.}
  \label{tab:components}
\end{table*}

\paragraph{Overview (\cref{tab:components}).} The table isolates one component per row. Multi-path targets and the translation block help on every benchmark where they apply, at no or negligible inference cost. Orbit aggregation helps on T-LESS, SYMSOL-T and SPEED+ at about $4\text{--}5\times$ the inference cost, but hurts on YCB-V, so it is an option selected on validation. Geometry fusion is the most application-dependent component: it gives the largest single gain on the nearly symmetric spacecraft, a small gain on T-LESS, and none on SYMSOL-T and YCB-V, at about $4\times$ the cost. Distillation is application-dependent as well: it improves the spacecraft and T-LESS systems in every setting (T-LESS validation $10.82^\circ\to10.08^\circ$ single-view, $8.28^\circ\to7.88^\circ$ with orbit aggregation), but on SYMSOL-T it gains only $0.02^\circ$ in the single view and makes the orbit system worse, so it is an option chosen on validation rather than a general improvement. For the shared architecture this means that the core (frozen DINOv3, hierarchical head, multi-path targets, translation block) transfers across applications, while the expensive options have to be chosen per application.

\begin{table}[t]
  \centering\footnotesize
  \setlength{\tabcolsep}{2.4pt}
  \resizebox{\linewidth}{!}{\begin{tabular}{@{}l c cc cc@{}}
\toprule
 & & \multicolumn{2}{c}{T-LESS (7 obj.)} & \multicolumn{2}{c}{SYMSOL-T (5)} \\
\cmidrule(lr){3-4}\cmidrule(lr){5-6} Training target & Inference & val & test & val & test \\
\midrule
single annotated label & 1 view & 12.3\,{\scriptsize$\pm$0.1} & 17.9\,{\scriptsize$\pm$0.2} & 3.87\,{\scriptsize$\pm$0.03} & 3.88\,{\scriptsize$\pm$0.01} \\
 & orbit (4) & 9.3\,{\scriptsize$\pm$0.1} & 16.6\,{\scriptsize$\pm$0.2} & 2.65\,{\scriptsize$\pm$0.01} & 2.67\,{\scriptsize$\pm$0.00} \\
\addlinespace[2pt]
canonical label & 1 view & 11.5\,{\scriptsize$\pm$0.0} & 16.7\,{\scriptsize$\pm$0.1} & 3.40\,{\scriptsize$\pm$0.01} & 3.39\,{\scriptsize$\pm$0.01} \\
 & orbit (4) & 10.7\,{\scriptsize$\pm$0.0} & 18.2\,{\scriptsize$\pm$0.7} & 9.73\,{\scriptsize$\pm$0.05} & 9.69\,{\scriptsize$\pm$0.05} \\
\addlinespace[2pt]
all equivalents (ours) & 1 view & 10.8\,{\scriptsize$\pm$0.1} & 15.7\,{\scriptsize$\pm$0.2} & 2.72\,{\scriptsize$\pm$0.01} & 2.71\,{\scriptsize$\pm$0.01} \\
 & orbit (4) & 8.3\,{\scriptsize$\pm$0.0} & 15.2\,{\scriptsize$\pm$0.3} & 2.10\,{\scriptsize$\pm$0.01} & 2.11\,{\scriptsize$\pm$0.01} \\
\addlinespace[2pt]
\;+ orbit distillation & 1 view & 10.1\,{\scriptsize$\pm$0.1} & 15.2\,{\scriptsize$\pm$0.5} & 2.63\,{\scriptsize$\pm$0.01} & 2.62\,{\scriptsize$\pm$0.01} \\
 & orbit (4) & 7.9\,{\scriptsize$\pm$0.1} & 14.5\,{\scriptsize$\pm$0.3} & 2.26\,{\scriptsize$\pm$0.00} & 2.27\,{\scriptsize$\pm$0.00} \\
\bottomrule
\end{tabular}}
  \caption{\textbf{Training targets for symmetric objects} (symmetry-aware mean error, PoE). Single annotated label (A0), canonical label closest to identity (A1), and all equivalent rotations (A2, \cref{eq:symloss}); the last block adds orbit distillation on top of the selected target.}
  \label{tab:sym}
\end{table}

\paragraph{Symmetry-aware targets.}
\label{sec:abl_sym}
\Cref{tab:sym} compares the three training targets. Canonical labels (A1) remove the contradiction between equivalent labels and improve single-view accuracy on both datasets (SYMSOL-T $3.87^\circ\to3.40^\circ$, T-LESS validation $12.3^\circ\to11.5^\circ$). However, they break orbit voting: a canonical label is not related across in-plane rotations by \cref{eq:roll}, so the views disagree and the SYMSOL-T orbit error rises to $9.7^\circ$. Multi-path targets (A2) keep the relation and are the best single-view and orbit models on SYMSOL-T ($2.72^\circ$ and $2.10^\circ$; pre-registered criterion passes by about $100\times$ the seed std).

\begin{table}[t]
  \centering\small
  \setlength{\tabcolsep}{3pt}
  \resizebox{\linewidth}{!}{\begin{tabular}{@{}llc l ccc@{}}
\toprule
Head (params/branch) & Grid & Floor & Inference & Synth. & LB & SL \\
\midrule
Tesseract tree (2.4\,M) & 131,072 leaves & 2.51 & 1 view & 4.85\,{\scriptsize$\pm$0.03} & 17.33\,{\scriptsize$\pm$0.31} & 20.15\,{\scriptsize$\pm$0.57} \\
 & &  & orbit & 3.87\,{\scriptsize$\pm$0.03} & 12.30\,{\scriptsize$\pm$0.50} & 12.25\,{\scriptsize$\pm$0.37} \\
Flat (36.1\,M) & Hopf 10,923$\times$12 & 6.06 & 1 view & 8.95\,{\scriptsize$\pm$0.06} & 20.56\,{\scriptsize$\pm$0.70} & 27.16\,{\scriptsize$\pm$1.04} \\
 & &  & orbit & 7.53\,{\scriptsize$\pm$0.01} & 13.15\,{\scriptsize$\pm$0.37} & 16.31\,{\scriptsize$\pm$0.59} \\
Flat, same grid & 131,072 leaves & 2.51 & 1 view & 8.19\,{\scriptsize$\pm$0.15} & 23.83\,{\scriptsize$\pm$0.44} & 32.87\,{\scriptsize$\pm$0.29} \\
 & &  & orbit & 5.34\,{\scriptsize$\pm$0.02} & 13.04\,{\scriptsize$\pm$0.32} & 16.73\,{\scriptsize$\pm$0.56} \\
Flat, balanced Hopf & Hopf 1,771$\times$74 & 2.29 & 1 view & 8.14\,{\scriptsize$\pm$0.09} & 24.05\,{\scriptsize$\pm$0.61} & 31.21\,{\scriptsize$\pm$0.60} \\
 & &  & orbit & 5.28\,{\scriptsize$\pm$0.02} & 12.88\,{\scriptsize$\pm$0.70} & 16.29\,{\scriptsize$\pm$0.29} \\
\bottomrule
\end{tabular}}
  \caption{\textbf{Tree vs.\ flat grid} (PoE, mean error in $^\circ$). Same frozen features, augmentation, training budget and orbit inference (4 views, top-8 per view). Floor: mean error of the nearest anchor (tree: its leaf) on the synthetic-validation rotations. The \emph{same-grid} flat head classifies the tree's own 131,072 leaves, isolating hierarchical prediction; the balanced Hopf grid tests an alternative discretisation of the same size. 3 seeds (± seed std).}
  \label{tab:flat}
\end{table}

\paragraph{Hierarchical tree vs.\ flat grid (SPEED+).}
\label{sec:abl_flat}
We replace the tree with a count-matched flat classifier over 131,076 Hopf-fibration anchors~\cite{yershova2010hopf} in SPACE-HOP's configuration (10,923 directions $\times$ 12 rolls), keeping features, augmentation, training budget and orbit inference identical (\cref{tab:flat}). The tree is far more accurate ($3.87^\circ$ vs.\ $7.53^\circ$ synthetic with orbit voting), better on real data in all comparisons, and uses $15\times$ fewer parameters per branch (2.4\,M vs.\ 36.1\,M, of which 33.6\,M is the 131k-way output layer). Orbit voting helps both heads, and the flat head with orbit voting even has slightly fewer severe flips on lightbox, so flip repair comes mainly from orbit voting.

That comparison alone does not isolate hierarchical prediction, because with only 12 in-plane rolls the SPACE-HOP grid has a quantisation floor of $6.06^\circ$ on the synthetic-validation rotations, against $2.51^\circ$ for the tree's leaves. We therefore trained two further flat heads with the identical recipe and evaluated them with identical orbit inference (4 views, top-8 candidates per view, orbit-product scoring): one classifies the tree's \emph{own} 131,072 leaf centres and scores an orbit candidate by the probability of its Tesseract cell, exactly as the tree does, so that only the prediction (hierarchical vs.\ flat) differs; the other classifies a balanced Hopf grid of 1,771 Fibonacci directions $\times$ 74 rolls (131,054 anchors), whose floor ($2.29^\circ$ mean, $3.8^\circ$ max) is \emph{lower} than the tree's. On the same grid, the flat head reaches $5.34^\circ$ with orbit aggregation and $8.19^\circ$ from one view, against $3.87^\circ$ and $4.85^\circ$ for the tree (seed std $\le0.15^\circ$; pre-registered criterion passes). The better balanced grid helps the flat head only marginally ($5.28^\circ$ and $8.14^\circ$). The same ordering holds for the DINOv3-only variant, on the augmented validation views and in mean error on both real domains (\cref{tab:flat}); with orbit aggregation the flat heads have slightly fewer severe flips on lightbox (2.3\% vs.\ 3.0\%) but a higher median ($6.3^\circ$ vs.\ $5.3^\circ$). The tree's advantage therefore comes from hierarchical prediction itself, not from its discretisation: a single 131k-way softmax trained with one example per pose spreads the data over many rarely seen classes, whereas the tree shares the statistics of every coarse level across all of its leaves, with about $14\times$ fewer parameters.

\paragraph{Cost of the output layer.} We timed the three heads in one interleaved run under identical conditions (seed-0 checkpoints, same crops, one shared GPU; upper bounds). The tree is the slower head on a GPU: at batch 1 its autoregressive decoding takes 34.3\,ms for the DINOv3 branch, against 9.7\,ms for either flat head, whose 131k-way output is a single matrix product, so the single-view pipeline runs in 58.5\,ms against 33.5\,ms; with orbit aggregation the head costs 184\,ms against 62--72\,ms, and at batch 64 0.53\,ms against 0.13\,ms per image. In exchange the tree has $15\times$ fewer parameters (2.39\,M vs.\ 36.1\,M per branch), trains about twice as fast per step (18 vs.\ 37\,ms at batch 256) with less memory (3.6 vs.\ 4.5\,GB), and is far more accurate; inference memory is dominated by the shared backbone and is similar (1.2 vs.\ 1.4\,GB). The tree's case is therefore accuracy, size and training cost, not GPU latency.

\begin{table}[t]
  \centering\small
  \setlength{\tabcolsep}{3pt}
  \begin{tabular}{@{}l ccc@{}}
\toprule
Selection over 32 candidates & Synth. & LB & SL \\
\midrule
1 view (greedy) & 4.85\,{\scriptsize$\pm$0.03} & 17.33\,{\scriptsize$\pm$0.31} & 20.15\,{\scriptsize$\pm$0.57} \\
max-confidence & 5.51\,{\scriptsize$\pm$0.03} & 14.63\,{\scriptsize$\pm$0.42} & 15.92\,{\scriptsize$\pm$0.15} \\
top-1 vote & 4.37\,{\scriptsize$\pm$0.03} & 13.24\,{\scriptsize$\pm$0.42} & 13.38\,{\scriptsize$\pm$0.57} \\
weighted vote (32) & 3.46\,{\scriptsize$\pm$0.03} & 12.58\,{\scriptsize$\pm$0.42} & 12.88\,{\scriptsize$\pm$0.49} \\
orbit product & 3.87\,{\scriptsize$\pm$0.03} & 12.30\,{\scriptsize$\pm$0.50} & 12.25\,{\scriptsize$\pm$0.37} \\
\bottomrule
\end{tabular}
  \caption{\textbf{Selection rule over the same 32 orbit candidates.}}
  \label{tab:scoring}
\end{table}

\paragraph{How to select among orbit candidates (SPEED+).}
With the same 32 candidates (\cref{tab:scoring}), picking the most confident one is clearly worse than any form of agreement, and voting among the four views' top-1 predictions is better. Using all candidates helps most. The kernel-weighted vote is best on synthetic data ($3.46^\circ$). The likelihood product of \cref{eq:orbit} is best on both real domains, but its advantage over the vote is within about one seed std. Its pre-registered criterion (``product beats voting on synthetic validation'') \emph{fails}, and under our selection protocol the weighted vote would be preferred. The product is reported in \cref{tab:main} only because it was fixed in the base recipe beforehand. We therefore claim multi-view consistency over hierarchical candidates, not a specific scoring rule.

\begin{table}[t]
  \centering\small
  \setlength{\tabcolsep}{2.5pt}
  \resizebox{\linewidth}{!}{\begin{tabular}{@{}l ccc ccc@{}}
\toprule
Augmentation & \multicolumn{3}{c}{1 view} & \multicolumn{3}{c}{orbit (4)} \\
\cmidrule(lr){2-4}\cmidrule(lr){5-7} & Syn. & LB & SL & Syn. & LB & SL \\
\midrule
none & 5.54\,{\scriptsize$\pm$0.07} & 26.90 & 36.66 & 4.62\,{\scriptsize$\pm$0.04} & 20.25 & 24.83 \\
photometric & 5.00\,{\scriptsize$\pm$0.03} & 17.38 & 19.68 & 4.25\,{\scriptsize$\pm$0.06} & 12.29 & 12.37 \\
in-plane roll & 4.91\,{\scriptsize$\pm$0.00} & 24.30 & 31.52 & 3.78\,{\scriptsize$\pm$0.03} & 18.69 & 20.70 \\
both & 4.85\,{\scriptsize$\pm$0.03} & 17.33 & 20.15 & 3.87\,{\scriptsize$\pm$0.03} & 12.30 & 12.25 \\
\bottomrule
\end{tabular}}
  \caption{\textbf{Which augmentation matters.} \emph{None} is the clean tree. The roll uses exact labels (\cref{eq:roll}).}
  \label{tab:aug}
\end{table}

\paragraph{Photometric vs.\ rotation augmentation (SPEED+).}
\label{sec:abl_aug}
\Cref{tab:aug} separates the two augmentation families. Photometric augmentation with orbit voting reaches $12.3^\circ/12.4^\circ$ on real images, the same as the combined recipe. Rotation augmentation alone gives the best synthetic precision ($3.78^\circ$) but leaves real errors at $18.7^\circ/20.7^\circ$. Orbit voting on the clean tree already removes a quarter to a third of the real error ($26.9^\circ\!\to\!20.3^\circ$, $36.7^\circ\!\to\!24.8^\circ$). The real-domain gain therefore comes from \emph{lighting robustness plus orbit consistency}. The exact roll transform is needed for orbit voting itself; as a training augmentation it mainly improves synthetic precision. This corrects our initial hypothesis that rotation augmentation causes the flip repair.


\paragraph{Can training learn orbit consistency (SPEED+).}
On SPEED+ we compared paired views of the same image (a data control), a transported consistency loss on the leaf distribution (a standard equivariance regulariser), a \emph{hierarchical} transported loss at every tree level, and orbit distillation (\cref{tab:consistency}, supplementary). The hierarchical loss beats the base recipe and the paired-view control but not the leaf-only loss; its pre-registered criterion fails, so we make no hierarchical-transport claim. Both consistency losses mostly improve synthetic precision. Orbit distillation is the only objective that clearly improves the real domains ($11.2^\circ/11.1^\circ$ with orbit voting).

\paragraph{Further ablations} in the supplementary material cover the decoder lineage (MLP, GRU and Transformer decoders, hard vs.\ soft targets, beam search), fusion (single branches, learned fusion, weighted PoE), finer tokens, continuous flow posteriors and verifiers, discrete flow matching, oracle headroom, zero-shot transfer and timing.
